\section{Energy comparisons and joint spatial bounds}\label{sp:section}

We need a supply of particles in most unit cells, and a small probability
that many prescribed cells contain too many close neighbors.  The estimates
must hold jointly for arbitrary sets of cells.  They will follow from local
energy penalties, tested against products of count cutoffs.

Throughout this section the torus is $\Lambda_K=(\R/K\Z)^3$, with integer
$K\ge100$, and
\[
 V=K^3,\qquad D=N/V,\qquad r=\alpha/D,
 \qquad \alpha_0/2\le\alpha\le2\alpha_0.
\]
The fixed constant $\alpha_0$ will be chosen large.  We then take $D$
sufficiently large.  Let $\mathcal L=(\Z/K\Z)^3$ and
$B_v=v+[-1/2,1/2)^3$, using their periodic images.  Distances below are torus
distances.  We write $E_n$ for the full labelled $n$-particle Dirichlet
ground energy on this same torus with core $r$, and $E_0=0$.
The function $\Phi$ is a normalized nonnegative symmetric ground function,
and $\mu(\dd X)=\Phi(X)^2\dd X$.

For a subset $I$ of the particle labels, put
\[
 q_I[F]=\sum_{i\in I}\int |\nabla_iF|^2,
 \qquad n_v^I(X)=\#\{i\in I:x_i\in B_v\},\qquad n_v=n_v^{\{1,\ldots,N\}}.
\]
Integrals of form functions use their zero extensions.  Constants $c,C$
are positive and may change; dependence on fixed cutoff profiles and on
the fixed window for $\alpha$ is allowed.  They never depend on $K,N,D$ or
the choice of ground function.  Any additional dependence is stated.

\subsection{Two local kinetic inequalities}

We use the hard-core two-particle coercivity underlying local exclusion
bounds \cite[Proposition~10]{LundholmPortmannSolovej2015} and the Neumann
subdivision method \cite{LiebYngvason1998}. The following ray and
subdivision argument proves the occupation estimate needed here.

\begin{lemma}[Cube energy]\label{sp:cube}
Let $Q$ be a Euclidean cube of side $l$, and let $F\in H^1(Q^n)$ vanish
almost everywhere whenever $|x_i-x_j|\le r$ for some $i\ne j$.
With no boundary condition on the outer faces,
\begin{equation}\label{sp:cube-bound}
 \sum_{i=1}^n\int_{Q^n}|\nabla_iF|^2
 \ge c_{\rm cube}\frac r{l^3}n(n-1)\int_{Q^n}|F|^2,
 \qquad c_{\rm cube}=\frac1{64\sqrt3}.
\end{equation}
Consequently, for every torus form function $F\in H_0^1(\Omega_N)$
and every fixed label subset $I\subset\{1,\ldots,N\}$,
\begin{equation}\label{sp:cell-energy}
 q_I[F]\ge c_{\rm cube}r\int\sum_{v\in\mathcal L}
                    n_v^I(n_v^I-1)|F|^2.
\end{equation}
For $D\ge2$ this implies
\begin{equation}\label{sp:energy-lower}
 \frac{E_N}{N}\ge c_E\alpha,
 \qquad c_E=c_{\rm cube}/2.
\end{equation}
\end{lemma}

\begin{proof}
Fix a partner $y\in Q$.  A first-coordinate slice $f$ vanishes where
$|x-y|\le r$.  On a ray $x=y+s\omega$ in the cube, let $R(\omega)$
be its terminal radius.  Convexity makes the ray intersection an interval,
and $R(\omega)\le\sqrt3l$.  For $r<s<R(\omega)$,
\[
 |f(y+s\omega)|^2
 \le\left(\int_r^s t^{-2}\dd t\right)
       \left(\int_r^s|\partial_tf(y+t\omega)|^2t^2\dd t\right)
 \le\frac1r\int_r^{R(\omega)}|\partial_tf|^2t^2\dd t.
\]
Integrating in $s$ and $\omega$ proves
\begin{equation}\label{sp:ray}
 \int_Q|f|^2\le\frac{\sqrt3l^3}{r}\int_Q|\nabla f|^2.
\end{equation}
Rays ending before $r$ contribute zero.  The argument first applies to
smooth restrictions and then to Sobolev slices by approximation.  In a
torus cube lift, Euclidean distance at most $r$ also implies torus distance
at most $r$, so the required zero region is present.

For $n\ge2$, subdivide $Q$ into $j^3$ equal cubes, where
$j=\lfloor(n/2)^{1/3}\rfloor$.  Then
$n/16\le j^3\le n/2$.  In every sector assigning the $n$ coordinates
to these cubes, at most $j^3$ coordinates are in singleton cubes.
For each of the at least $n/2$ remaining coordinates choose a partner in
its cube, using a fixed label ordering.  Applying \eqref{sp:ray} to each
such coordinate gives at least
\[
 \frac{rj^3}{\sqrt3l^3}\frac n2
 \ge\frac{rn^2}{32\sqrt3l^3}
\]
times the sector norm.  We only restrict the coordinate integrals to
sectors; no derivative of a sector indicator is taken.  Sum these
inequalities over sectors.  The cases $n=0,1$ have zero lower bound,
and impossible sectors have zero norm.  This proves
\eqref{sp:cube-bound}, with the stated smaller constant.

Apply the cube bound to all unit-cell assignment sectors of the labels
in $I$, fixing the other coordinates, and integrate them.  This gives
\eqref{sp:cell-energy}.  Finally,
$\sum_v n_v(n_v-1)\ge N^2/V-N$, so
$E_N/N\ge c_{\rm cube}r(D-1)\ge c_E\alpha$.
\end{proof}

\begin{lemma}[Close-neighbor energy]\label{sp:close-energy}
Let $0<r<s<1$, let $F\in H_0^1(\Omega_N)$, and fix a label set
$I\subset\{1,\ldots,N\}$.  Define
$d_i^I=\min_{j\in I\setminus\{i\}}d(x_i,x_j)$, with value $+\infty$
when the minimizing set is empty.  Then
\begin{equation}\label{sp:close-energy-bound}
 q_I[F]\ge c_{\rm near}\,rs^{-3}
        \int\#\{i\in I:d_i^I\le2s\}|F|^2,
 \qquad c_{\rm near}=\frac{27}{262144\sqrt3}.
\end{equation}
\end{lemma}

\begin{proof}
Partition each coordinate circle into $\lfloor K/(8s)\rfloor$ equal
intervals and translate the resulting grid uniformly.  Its cube side
$l$ belongs to $[8s,16s]$.  Two points at distance at most $2s$ lie in
one cube with probability at least $(1-2s/l)^3\ge27/64$.
For each coordinate having a partner in its cube, \eqref{sp:ray}
gives kinetic energy at least $r/(\sqrt3l^3)$ times its squared norm.
Choose a nearest partner for each $i$ with $d_i^I\le2s$, and average
the restricted inequalities over the grid translation.  Each such
coordinate is counted with probability at least $27/64$.  Since
$l^3\le4096s^3$, this proves the displayed constant.
\end{proof}

\subsection{Normalized insertion and energy differences}

We next compare energies at different particle numbers.  Normalizing the
inserted orbital separately for each old configuration prevents the old
kinetic energy from acquiring a multiplicative error. The trial follows the
nearest-neighbor construction of Lieb, Seiringer, and Yngvason
\cite[Section~3.2]{LiebSeiringerYngvason2000}; the configurationwise
normalization below preserves the old marginal for arbitrary old form
functions.

\begin{lemma}[Normalized insertion]\label{sp:insertion}
Let $w$ be a real normalized smooth one-particle orbital and $F$ an
$n$-particle form function.  For every configuration $X$ in the support
of $F$, suppose at most $k$ old positions lie within distance $2r$ of
$\supp w$.  If
\begin{equation}\label{sp:insertion-smallness}
 C\norm{w}_\infty^2kr^3\le\frac14,
\end{equation}
there is an admissible function $\widetilde F$ with one more coordinate,
having the same norm as $F$, such that
\begin{equation}\label{sp:insertion-cost}
 q_{n+1}[\widetilde F]\le q_n[F]
       +C\left(\norm{\nabla w}_2^2+\norm{w}_\infty^2kr\right)\norm F_2^2.
\end{equation}
The constants depend only on a fixed radial cutoff profile.
\end{lemma}

\begin{proof}
Choose a nondecreasing Lipschitz function $h_r$ between zero and one, zero up to
$1.1r$, one beyond $2r$, and with Lipschitz constant at most $C/r$.
For old positions $X$ put
\[
 f_X(y)=w(y)\min_{1\le i\le n}h_r(d(y,x_i)),\qquad
 Z(X)=\int f_X(y)^2\dd y.
\]
The empty minimum is one.  The union of the relevant radius-$2r$ balls
has volume at most $Ckr^3$, so $Z\ge3/4$ under
\eqref{sp:insertion-smallness}.  Where a nearest minimizer is unique,
only that old coordinate and $y$ contribute to its cutoff derivative.
The same bound holds almost everywhere for the Lipschitz minimum.
Therefore
\begin{equation}\label{sp:unnormalized-gradient}
 \int\left(|\nabla_yf_X|^2+
              \sum_{i=1}^n|\nabla_if_X|^2\right)\dd y
 \le C\left(\norm{\nabla w}_2^2+\norm w_\infty^2kr\right).
\end{equation}
Indeed cutoff derivatives have squared size at most $C/r^2$ and
are supported in the relevant balls.

Set $a_X=f_X/\sqrt Z$ wherever $Z>1/4$.  For an old-coordinate
derivative, normalization gives the orthogonal-projection identity
\[
 \partial a_X=Z^{-1/2}
       \bigl(\partial f_X-a_X\langle a_X,\partial f_X\rangle_{L^2(\dd y)}\bigr).
\]
It follows that its squared $L^2(\dd y)$ norm is at most
$Z^{-1}\norm{\partial f_X}_2^2$.
The new-coordinate derivative has the same upper bound because $Z$
does not depend on $y$.  Also
\[
 \int a_X^2\dd y=1,
 \qquad\int a_X\nabla_i a_X\dd y=0.
\]
Thus the cross term vanishes on differentiating $F(X)a_X(y)$ and
integrating $y$.  Equation~\eqref{sp:unnormalized-gradient} proves
\eqref{sp:insertion-cost}.

For a globally defined multiplier replace $\sqrt Z$ by
$\max(\sqrt Z,1/2)$.  The resulting $y$-norm is at most one and
equals one on a neighborhood of the support where the hypothesis is
used.  A Sobolev function and its weak gradient vanish almost everywhere
on its zero set, so this extension changes none of the preceding integrals.
The extra clearance $1.1r$ and compactly supported approximation of the
old function show that $\widetilde F=Fa_X$ belongs to the enlarged
Dirichlet form domain.  This also justifies the calculation for arbitrary
old form functions.
\end{proof}

\begin{corollary}[Energy differences]\label{sp:removal}
For all sufficiently large $D$, uniformly in the stated $\alpha$ window,
\begin{equation}\label{sp:removal-bound}
 E_N-E_{N-m}\le C\alpha m\qquad(0\le m\le N).
\end{equation}
Moreover $E_n/n$ is nondecreasing among positive particle numbers with
nonempty form domain, at fixed $K,r$.
\end{corollary}

\begin{proof}
Use $w=V^{-1/2}$ in Lemma~\ref{sp:insertion}.  For at most $N$ old
particles its norm loss is at most $CDr^3=C\alpha^3/D^2$, and its
energy cost is at most $CDr=C\alpha$.  Starting from the constant
one-particle function constructs nonzero trials at every particle number
up to $N$.  Apply the variational principle at each addition and sum
the costs to obtain \eqref{sp:removal-bound}.

For the last assertion, fix $k<n$ and average the kinetic energy of
all $k$-label subsets of a normalized $n$-particle form function.  Each
coordinate is counted with probability $k/n$, whereas slicing in the
other coordinates bounds each subset energy below by $E_k$.
Consequently $(k/n)q_n[F]\ge E_k$; taking the infimum gives
$E_n/n\ge E_k/k$.
\end{proof}

\subsection{A joint bound for deficient cells}

We apply the ground-state multiplier identity,
Equation~\eqref{setup:multiplier}, to products of local count cutoffs.

\begin{proposition}[Joint low-count bound]\label{sp:low-counts}
One can fix $\alpha_0$ sufficiently large and then $\ell>0$ sufficiently
small so that, for every $A\subset\mathcal L$ and all sufficiently large
$D$,
\begin{equation}\label{sp:low-count-bound}
 \mu\{n_v<\ell D\text{ for every }v\in A\}
       \le (C/D^2)^{|A|}.
\end{equation}
These choices and the constant $C$ are independent of $K$ and of the
ground function.
\end{proposition}

\begin{proof}
Fix translates $w_v$ of a smooth normalized real orbital supported
strictly inside $B_v$.  Choose translates $\chi_v\in[0,1]$ supported
inside $B_v$, equal to one on a fixed neighborhood of $\supp w_v$,
and satisfying $|\nabla\chi_v|^2\le C\chi_v$.  Put
$m_v=\sum_i\chi_v(x_i)$.  The neighborhood is chosen before $r$,
so for sufficiently small $r$ any old center whose cutoff ball reaches
$\supp w_v$ has $\chi_v=1$.

Suppose $F$ vanishes unless $m_v\le2\ell D$ for all $v\in A$.
The case $F=0$ is trivial, so assume $\norm F_2>0$.
Insert $\lfloor\ell D\rfloor$ particles at each of these sites,
successively using $w_v$.  At each insertion at most $3\ell D$
old centers can affect its orbital.  Supports at distinct sites have a
fixed positive separation, so insertions at other sites cause no additional
cutoff.  Lemma~\ref{sp:insertion} applies for large $D$ and gives
cost at most $C(1+3\ell\alpha)$ per inserted particle.

First choose $\alpha_0$ large, and then $\ell$ small, so that this
cost is at most $c_E\alpha/2$ throughout the $\alpha$ window.
Writing $k=\lfloor\ell D\rfloor|A|$, the enlarged trial has norm
$\norm F_2$ and energy at most $q_N[F]+(c_E\alpha/2)k\norm F_2^2$.
In particular it is nonzero, so $E_{N+k}<\infty$.
The monotonicity in Corollary~\ref{sp:removal} and
\eqref{sp:energy-lower} give
$E_{N+k}\ge E_N+kE_N/N\ge E_N+c_E\alpha k$.
The variational principle therefore proves
\begin{equation}\label{sp:deficiency-penalty}
 q_N[F]-E_N\norm F_2^2
       \ge c\alpha D|A|\norm F_2^2,
\end{equation}
where $c>0$ includes the fixed $\ell$ and $D$ is large enough that
$\lfloor\ell D\rfloor\ge\ell D/2$.

Let $\zeta:[0,\infty)\to[0,1]$ be a Lipschitz function equal to
one up to $1$ and zero from $2$ onwards, and set
$u_A=\prod_{v\in A}\zeta(m_v/(\ell D))$.
Since distinct $\chi_v$ have disjoint supports,
\[
 |\nabla m_v|^2=\sum_i|\nabla\chi_v(x_i)|^2\le Cm_v,
 \qquad
 |\nabla u_A|^2\le\frac C D
               \sum_{v\in A}u_{A\setminus\{v\}}^2.
\]
The second estimate uses $m_v\le2\ell D$ wherever the corresponding
cutoff derivative is nonzero.  Apply \eqref{sp:deficiency-penalty}
and \eqref{setup:multiplier} to $u_A\Phi$.  With
$a(A)=\int u_A^2\dd\mu$, this gives
\[
 a(A)\le\frac{C}{D^2|A|}\sum_{v\in A}a(A\setminus\{v\})
 \quad(A\ne\varnothing).
\]
Starting from $a(\varnothing)=1$, induction on $|A|$ yields
$a(A)\le(C/D^2)^{|A|}$.  Finally $m_v\le n_v$, so every
configuration in the event of \eqref{sp:low-count-bound} has $u_A=1$.
\end{proof}

The same constants work for arbitrarily many prescribed cells.  We now
obtain an upper count bound and a close-neighbor bound with that same joint
scope.  Their common input is a localized energy inequality.

\subsection{Exponential spatial localization}

For nonempty $A\subset\mathcal L$, use its site centers as a subset of
the torus and define
\begin{align}
 f_A(X)&=\sum_{i=1}^N e^{-d(x_i,A)},\label{sp:weighted-count}\\
 k_A(X)&=\#\{i:d(x_i,A)\le6,\ d_i\le2s_0\},
 \qquad d_i=\min_{j\ne i}d(x_i,x_j),\label{sp:localized-close-count}\\
 s_0&=bD^{-1/3},\qquad 0<b<1.\label{sp:separation-scale}
\end{align}
The parameter $b$ will be fixed small below.  We use $D$ large enough
that $r<s_0<1/2$.

The particle-number localization uses the many-body IMS construction;
compare \cite[Section~3]{Lewin2011}. We prove the product partition and its
weighted penalty explicitly.  We separate particles near $A$ from the
remaining particles.  The latter retain at least their own ground energy;
the former have enough local kinetic energy to pay the removal cost and
leave a positive penalty for excessive counts or close neighbors.

\begin{lemma}[Weighted local penalty]\label{sp:localized-form}
There are constants $C,c>0$, independent of $b$, such that every
$N$-particle form function $F$ satisfies
\begin{equation}\label{sp:localized-form-bound}
 q_N[F]-E_N\norm F_2^2
   \ge\alpha\int\bigl(f_A-CD|A|+cb^{-3}k_A\bigr)|F|^2.
\end{equation}
\end{lemma}

\begin{proof}
\emph{Partitioning the particles.}
Let $H\ge0$ be an auxiliary random radius with
$\PP(H\ge h)=e^{-h}$.  Choose a fixed Lipschitz profile
$\vartheta$ equal to $\pi/2$ on $(-\infty,0]$, zero on $[1,\infty)$,
and with bounded derivative.  Set
\[
 a_H(x)=\sin\vartheta(d(x,A)-H),\qquad
 c_H(x)=\cos\vartheta(d(x,A)-H).
\]
Thus $a_H^2+c_H^2=1$, $a_H=1$ at distance at most $H$, and
its support is within distance $H+1$.  Both squared gradients sum to
at most $C\ind_{\{H<d(x,A)<H+1\}}$.
For $I\subset\{1,\ldots,N\}$ define
\[
 F_{I,H}(X)=F(X)\prod_{i\in I}a_H(x_i)
                         \prod_{j\notin I}c_H(x_j).
\]
Expanding the gradients, using $a_H\nabla a_H+c_H\nabla c_H=0$,
gives the exact localization identity
\begin{equation}\label{sp:ims}
 \sum_I q_N[F_{I,H}]
 =q_N[F]+\int\sum_i\bigl(|\nabla a_H(x_i)|^2+
                        |\nabla c_H(x_i)|^2\bigr)|F|^2.
\end{equation}
The squared norms sum to $\norm F_2^2$.

\smallskip\noindent\emph{The energy retained outside the localized set.}
The index set $I$ consists of the particles localized near $A$.
For each $I$, slicing in the coordinates of $I$ and dropping their
exclusions from the others gives
\[
 q_{I^c}[F_{I,H}]\ge E_{N-|I|}\norm{F_{I,H}}_2^2
                \ge(E_N-C\alpha|I|)\norm{F_{I,H}}_2^2.
\]
To sum these removal costs, regard the squared product factors at fixed
$X,H$ as probabilities of independent choices of membership in $I$,
with probabilities $a_H(x_i)^2$.  In particular,
\[
 \sum_I |I|\prod_{i\in I}a_H(x_i)^2
                 \prod_{j\notin I}c_H(x_j)^2
       =\sum_i a_H(x_i)^2.
\]
Since
\[
 \E_H a_H(x)^2\le e\,e^{-d(x,A)},\qquad
 \E_H\ind_{\{H<d(x,A)<H+1\}}\le e\,e^{-d(x,A)},
\]
the mean removal cost and the mean localization error together are
at most $C\alpha\int f_A|F|^2$; we used the fixed positive lower
bound for $\alpha$ to absorb the error without $\alpha$.

\smallskip\noindent\emph{The local energy that pays for removal.}
The kinetic energy in the coordinates of $I$ is still available.  Apply half of
\eqref{sp:cell-energy} and half of \eqref{sp:close-energy-bound}
to $q_I[F_{I,H}]$.  Put $p_v=e^{-d(v,A)}$.
If $H\ge d(v,A)+1$, the whole centered unit cell $B_v$ is on the
$a_H=1$ side.  Summing over $I$ and averaging $H$ therefore makes
the first term at least
\[
 cr\int\sum_v p_v n_v(n_v-1)|F|^2.
\]
For a coordinate counted by $k_A$, some partner lies within $2s_0<1$,
and hence both are at distance less than $7$ from $A$.
On $H\ge8$ both labels are certainly in $I$.  This event has fixed
probability $e^{-8}$, so the second kinetic term is at least
\[
 crs_0^{-3}\int k_A|F|^2
        =c\alpha b^{-3}\int k_A|F|^2.
\]
In particular the partner used here belongs to the removed set before
the subset-label inequality is applied.

Combining these estimates with \eqref{sp:ims} gives
\begin{equation}\label{sp:before-square}
 q_N[F]-E_N\norm F_2^2
 \ge\int\left(cr\sum_vp_v n_v(n_v-1)
                 -C\alpha f_A+c\alpha b^{-3}k_A\right)|F|^2.
\end{equation}
\smallskip\noindent\emph{Leaving a positive count penalty.}
Within each unit cell its distance to $A$ differs from its center's
distance by at most $1$.  Consequently
\[
 f_A\le e\sum_vp_v n_v,
 \qquad
 \sum_vp_v\le\sum_{w\in A}\sum_v e^{-d(v,w)}\le C|A|.
\]
The last lattice exponential sum is uniformly bounded on every periodic
unit grid.  Since $r=\alpha/D$, completing a scalar square gives,
for all integers $n\ge0$ and $D\ge1$,
\[
 \frac cD n(n-1)-(C+1)e n\ge-C'D.
\]
Multiply by $\alpha p_v$ and sum.  In \eqref{sp:before-square}
this leaves the positive term $\alpha f_A$ and proves
\eqref{sp:localized-form-bound}.  None of the constants in this
calculation depends on $b$.
\end{proof}

\subsection{Joint count and close-neighbor tails}

\begin{proposition}[Joint spatial bounds]\label{sp:joint-tails}
For every $A\subset\mathcal L$,
\begin{equation}\label{sp:count-mgf}
 \log\int e^{zf_A}\dd\mu\le CzD|A|\qquad(0\le z\le1),
\end{equation}
where $f_\varnothing=0$.
Fix any $\theta>0$.  One can choose $b>0$ sufficiently small depending
on $\theta$, then take $D$ sufficiently large, so that the random site set
\begin{equation}\label{sp:close-bad-set}
 \mathcal Q(X)=\left\{v\in\mathcal L:
       \#\{i:d(x_i,v)\le5,\ d_i\le s_0\}\ge\theta D\right\}
\end{equation}
satisfies
\begin{equation}\label{sp:close-tail}
 \mu\{A\subset\mathcal Q\}\le p_D^{|A|},\qquad
 p_D=\left(C_{\theta}b^{-2}D^{-4/3}\right)^{1/Q}\longrightarrow0.
\end{equation}
Here $Q$ is an absolute integer; $D$ is taken large enough that the
displayed base is at most one.  All statements are uniform for
$\alpha\in[\alpha_0/2,2\alpha_0]$, $K\ge100$, and the ground function.
\end{proposition}

\begin{proof}
Distance is $1$-Lipschitz and $e^{-d}\le1$, so
$|\nabla f_A|^2\le f_A$ almost everywhere.  Let
$M(z)=\int e^{zf_A}\dd\mu$.  Apply
\eqref{setup:multiplier} and \eqref{sp:localized-form-bound}
with $u=e^{zf_A/2}$, dropping the nonnegative close-count term.
This gives
\[
 \alpha\bigl(M'(z)-CD|A|M(z)\bigr)
       \le\frac{z^2}{4}M'(z).
\]
Our choice of $\alpha_0$ ensures $\alpha\ge1$.  For $0\le z\le1$
we therefore obtain $M'(z)/M(z)\le CD|A|$.  Integrating from
$M(0)=1$ proves \eqref{sp:count-mgf}.  These differentiations are
ordinary finite-volume ones, since $0\le f_A\le N$.

For the second statement first assume that distinct sites of $A$ have
distance greater than $20$.  Choose translates $a_v\in[0,1]$, equal
to one on the radius-$5$ ball and supported on the radius-$6$ ball,
with $|\nabla a_v|^2\le Ca_v$.  Let $\chi\in[0,1]$ equal one on
$[0,1]$, zero on $[2,\infty)$, and satisfy $|\chi'|^2\le C\chi$.
Such cutoffs are squares of fixed Lipschitz profiles.  Define
\[
 \widetilde m_v(X)=\sum_i a_v(x_i)\chi(d_i/s_0).
\]
Then $\sum_{v\in A}\widetilde m_v\le k_A$, and each close count
in \eqref{sp:close-bad-set} is at most $\widetilde m_v$.

We verify the gradient estimate
\begin{equation}\label{sp:nearest-gradient}
 |\nabla\widetilde m_v|^2\le Cs_0^{-2}\widetilde m_v.
\end{equation}
Away from a null set, every nearest neighbor is unique.  If $x_i,x_j$
have the same nearest neighbor $x_k$, with the first incoming distance
at least the second, then $|x_i-x_j|\ge|x_i-x_k|$.
The cosine law shows that the two directions at $x_k$ have angle at
least $\pi/3$.  For the edges of length at most $2s_0$ used by the
derivatives, consistent Euclidean lifts exist because $s_0<1/2$ and
$K\ge100$.  Packing directions on the unit sphere bounds the incoming
degree by a numerical constant.  Each coordinate derivative therefore
meets only a bounded number of summands of $\widetilde m_v$.
The inequalities for $a_v,\chi$ bound the squared full gradient of
each summand by $Cs_0^{-2}a_v(x_i)\chi(d_i/s_0)$.
Summing with bounded overlap proves \eqref{sp:nearest-gradient}.
All coordinates in those derivatives are within $6+2s_0<7$ of $v$.

Choose a Lipschitz increasing function $\eta$ equal to zero on
$[0,1/2]$ and one on $[1,\infty)$, and put
\[
 u_A=\prod_{v\in A}\eta(\widetilde m_v/(\theta D)),
 \qquad a(A)=\int u_A^2\dd\mu.
\]
On the nonzero set of $u_A$ we have
$k_A\ge(\theta D/2)|A|$.  Choose $b$ so small that the constants
in \eqref{sp:localized-form-bound} satisfy
$cb^{-3}\theta/2-C\ge1$.  That form inequality gives
\[
 q_N[u_A\Phi]-E_N\norm{u_A\Phi}_2^2
       \ge\alpha D|A|a(A).
\]
At almost every configuration, gradients belonging to distinct sites involve
disjoint particle coordinates, by the separation of the sites.
Using \eqref{sp:nearest-gradient} on the transition region
$\widetilde m_v\le\theta D$ therefore gives
\[
 |\nabla u_A|^2\le\frac{C}{\theta Ds_0^2}
                    \sum_{v\in A}u_{A\setminus\{v\}}^2.
\]
Equation~\eqref{setup:multiplier} and induction from
$a(\varnothing)=1$ now imply
\[
 a(A)\le\left(\frac{C}{\alpha\theta D^2s_0^2}\right)^{|A|}
       \le\left(C_\theta b^{-2}D^{-4/3}\right)^{|A|}.
\]
On $A\subset\mathcal Q$, every cutoff equals one, proving this
bound for separated $A$.

Finally the graph joining distinct lattice sites at distance at most
$20$ has uniformly bounded degree, also on the torus.  It admits a
coloring with an absolute number $Q$ of colors.  Every prescribed
$A$ contains a separated color class $A'$ of size at least $|A|/Q$.
The event $A\subset\mathcal Q$ implies $A'\subset\mathcal Q$.
Applying the separated estimate, with its base at most one, proves
\eqref{sp:close-tail} for all $A$.
\end{proof}

These estimates concern the actual law $\mu$, without conditioning on
the bath or on future path data.  A close-pair defect defined using only
the bath coordinates is contained in the corresponding full-configuration
defect.  It therefore inherits \eqref{sp:close-tail} under the bath
marginal by integrating out the first coordinate.  This inclusion will
be used for the one-coordinate geometry.

The choices have a fixed order: $\alpha_0$ large, then $\ell$ small;
next any required fixed $\theta$, then $b$ small; finally $D$ large.
Every lower threshold on $D$ above is independent of $K\ge100$.
In particular these spatial bounds remain uniform when $D$ varies in
a compact subinterval of its allowed range along a thermodynamic sequence.
